\section{What would prove this wrong}\label{sec:refute}

Each claim below is paired with the cheapest observation that would falsify it. A test that needs
a new instrument is not a test; none of these does.

\begin{center}
\begin{tabular}{@{}p{0.42\textwidth}p{0.52\textwidth}@{}}
\toprule
claim & what would refute it \\
\midrule
\S\ref{sec:redistribution}: $\Fr$ is traceless and $\Fr=0\iff[\Omega,\PP]=0$ &
a single matrix pair with nonzero trace, or with $\Fr=0$ and $[\Omega,\PP]\neq0$. None exists;
the statements are kernel-checked. Refutation here would mean a bug in Lean. \\
\S\ref{sec:geometry}: $g$ is positive semidefinite &
a projector family with $g(\dot\PP,\dot\PP)<0$ at some point. \\
\S\ref{sec:geometry}: $\abs{\operatorname{Im}Q}\le\sqrt{g_{\mu\mu}g_{\nu\nu}}$ &
a pair of parameters with the ratio exceeding $1$. \\
\S\ref{sec:hierarchy}: rank one gains nothing &
a family with $\norm{\dot P}_{\mathrm{op}}<\norm{\dot v_\perp}$ or
$\norm{\dot P}_F<\sqrt2\norm{\dot v_\perp}$. \\
\S\ref{sec:hierarchy}: cluster projector rate $\norm{\dot A}/\Gap_{\mathrm{ext}}$ &
a family whose cluster projector moves faster than $\norm{\dot A}/\Gap_{\mathrm{ext}}$ (would
refute the constant, not the principle). \\
\S\ref{sec:totals}: $\sum_{i<j}\norm{\nabla d_{ij}}^2=N(N-2)$ &
one configuration of $N$ equal masses with a different total. This is the cheapest test in the
paper: it is a finite computation for any $N$, and a single exception kills it. \\
\S\ref{sec:shape}: $\det M^*M=12A^2/(xyz)$ &
one triangle with a different determinant. Likewise one computation. \\
\S\ref{sec:shape}: the cone is isotropic and linear &
a measurement of the angular dependence of $\sigma_1^2-\sigma_2^2$ near an equilateral shape with
a slope other than $3$, or direction-dependent. \\
\S\ref{sec:shape}: holonomy $-1$ &
a loop around an equilateral shape returning the eigen-line to itself. \\
\S\ref{sec:refraction}: Snell for conformal metrics &
nothing: it is a theorem. What is refutable is the modelling claim that a given physical medium
is conformal with factor $\rho$, and that is refuted by measuring the index directly. \\
\bottomrule
\end{tabular}
\end{center}

\begin{remark}[The asymmetry this paper is trying to fix]
Nine of the ten papers in the corpus say what they explain. One says what would refute it. The
table above is an attempt to make the ratio one to one.
\end{remark}
